% ============================================================
% CAPITOLO 4 -- TEST AUTOMATICI
% ============================================================
\section{Test automatici}\label{sec:automatici}

Oltre ai test manuali basati su Category-Partition, sono stati
applicati tre approcci di generazione automatica dei test, ciascuno
basato su un principio diverso: il \emph{random testing}, che genera
input in modo casuale tramite reflection (Randoop); la generazione
assistita da un modello di intelligenza artificiale (LLM), a
partire da diverse formulazioni del prompt; la generazione
\emph{coverage-guided}, che ottimizza la ricerca di nuovi test verso
la massimizzazione della copertura del codice (EvoSuite). Per ciascuna
classe, i tre approcci sono stati applicati e confrontati tra loro,
evidenziando punti di forza e limiti complementari.
\subsection{Test automatici su \texttt{Resolver}}

\paragraph{Randoop}


Randoop (versione 4.3.3) è uno strumento di \emph{feedback-directed
random testing}: genera sequenze di chiamate a metodi con input
casuali, sfruttando la reflection Java. Lo strumento richiede di essere
eseguito con lo stesso JDK usato per compilare il progetto (JDK 11):
se lanciato con una JVM più recente, genera codice che utilizza API non
disponibili nella versione richiesta per la compilazione finale (ad
esempio \texttt{isEmpty()} su \texttt{CharSequence}, introdotto in
Java 15).

Con un budget di generazione di 60 secondi sono stati generati 3508
test, nessuno dei quali error-revealing. La \tabref{tab:randoop-results} 
riporta i risultati numerici completi
per entrambe le classi; le \tabref{fig:jacoco-resolver-randoop} e
\tabref{fig:jacoco-utf8-randoop} mostrano il dettaglio per metodo. La
copertura della sola suite Randoop, misurata isolatamente, risulta
molto bassa su \texttt{Resolver}. La causa è che la maggior parte dei
metodi coinvolti richiede come precondizione oggetti complessi
(\texttt{Schema}, \texttt{GenericData}) che Randoop non è in grado di
costruire autonomamente tramite reflection. I test generati seguono
quasi tutti lo stesso pattern: costruzione di uno \texttt{Schema}
vuoto o con argomenti \texttt{null}, invocazione di
\texttt{resolve()} e cattura dell'eccezione risultante (tipicamente
\texttt{NullPointerException}), raggiungendo il metodo solo al punto
di ingresso senza mai esercitarne la logica decisionale interna.

Dal punto di vista qualitativo, i test generati utilizzano nomi
sequenziali privi di significato semantico (ad esempio \texttt{test1},
\texttt{test2}), senza commenti né documentazione, con una struttura
lineare di chiamate priva di separazione logica: comprendere l'intento
di un singolo caso di test richiede la lettura dell'intera
implementazione.

\paragraph{Generazione tramite LLM}

Sono stati utilizzati diversi prompt per sperimentare la generazione
di test da parte di un LLM. L'IA utilizzata è stata Gemini,
somministrando 3 prompt con quantità di contenuto crescente.
Il testo completo dei tre prompt è riportato nella
\tabref{tab:prompt-resolver}; la \tabref{tab:llm-resolver-results} ne
riassume i risultati numerici, mentre la
\tabref{fig:jacoco-resolver-llm} mostra il dettaglio di copertura per
metodo della suite risultante (Prompt 2 e 3 combinati).

Il primo prompt, formulato senza fornire il codice sorgente della
classe, ha prodotto codice non compilabile: il modello ha introdotto
due classi di eccezione inesistenti nel progetto e un nome di classe
errato, presumibilmente per analogia con librerie di validazione più
comuni. Il secondo prompt, che includeva il codice sorgente completo,
non ha prodotto alcuna allucinazione sui nomi delle API utilizzate,
richiedendo solo l'appiattimento di una struttura basata su
\texttt{@Nested} (incompatibile con la configurazione Surefire del
progetto): dopo questa correzione, tutti i 45 test generati sono
risultati passanti. Il terzo prompt, che aggiungeva vincoli di
progetto espliciti (assenza di \texttt{@Nested}, stile di
formattazione, categorie di risoluzione da coprire sistematicamente),
ha prodotto 13 test integrabili senza alcuna correzione.

Il salto di qualità più rilevante si osserva tra il primo e il secondo
prompt: la disponibilità del codice sorgente elimina quasi
completamente le allucinazioni sui nomi delle API. Il passaggio dal
secondo al terzo prompt migliora invece principalmente
l'\emph{integrabilità} del codice generato, riducendo a zero gli
interventi manuali necessari, più che la correttezza di base del
codice prodotto.


\paragraph{EvoSuite}

EvoSuite è uno strumento per la generazione automatica di test guidata
dalla copertura, eseguito con JDK 11 per lo stesso motivo riportato in
precedenza per Randoop. La generazione è stata eseguita con un budget
di 60 secondi, producendo 78 test. La \tabref{tab:evosuite-resolver}
riassume i risultati, e la \tabref{fig:jacoco-resolver-evosuite} ne
mostra il dettaglio per metodo.

Durante la generazione, EvoSuite ha segnalato un test instabile
(\emph{flaky}): un valore non deterministico utilizzato internamente
nel nome di uno schema non risolto varia tra un'esecuzione e l'altra.

La copertura, misurata isolatamente con JaCoCo, risulta pari
all'88.5\% statement e all'84.2\% branch — la più alta tra le tre
tecniche automatiche applicate a \texttt{Resolver}, risultato atteso
poiché EvoSuite è l'unica delle tre esplicitamente guidata dal
criterio di copertura durante la generazione.

\subsection{Test automatici su \texttt{Utf8}}

\paragraph{Randoop}

Randoop è stato eseguito su \texttt{Utf8} con la stessa configurazione
già descritta per \texttt{Resolver} (versione 4.3.3, JDK 11, budget di
60 secondi), generando 2841 test, nessuno dei quali error-revealing.
La \tabref{tab:randoop-results} riporta i risultati numerici completi
per entrambe le classi; la \tabref{fig:jacoco-utf8-randoop} mostra il
dettaglio per metodo.

A differenza di \texttt{Resolver}, su \texttt{Utf8} Randoop ottiene una
copertura molto più alta (79\% statement, 75\% branch). Il motivo è
speculare a quanto osservato in precedenza: \texttt{Utf8} accetta
input primitivi semplici (\texttt{String}, \texttt{byte[]}, interi),
che Randoop riesce a costruire facilmente tramite reflection, senza
dover generare oggetti compositi complessi come \texttt{Schema} e
\texttt{GenericData}. Questo contrasto conferma empiricamente che il
limite riscontrato su \texttt{Resolver} non è generale allo strumento,
ma specifico ai domini con oggetti compositi articolati.
\paragraph{Generazione tramite LLM}

Sono stati sperimentati gli stessi tre tipi di interrogazioni già
utilizzati per \texttt{Resolver}, con Gemini come modello. Il testo
completo dei tre prompt è riportato nella \tabref{tab:prompt-utf8}; la
\tabref{tab:llm-utf8-results} ne riassume i risultati numerici, mentre
la \tabref{fig:jacoco-utf8-llm} mostra il dettaglio di copertura per
metodo.

A differenza di \texttt{Resolver}, il primo prompt (senza codice
sorgente) non ha prodotto allucinazioni evidenti sui nomi dei metodi:
un'ipotesi plausibile è che \texttt{Utf8}, essendo una classe
pubblica ampiamente documentata e utilizzata negli esempi Avro, sia
più "nota" al modello dal materiale di addestramento rispetto a una
classe più interna come \texttt{Resolver}. Il secondo prompt ha
prodotto 23 test, tutti passanti dopo l'appiattimento della struttura
\texttt{@Nested}. Il terzo prompt ha prodotto 17 test, di cui 16
integrabili senza correzioni: un singolo test è stato rimosso perché
utilizzava un costruttore package-private non accessibile dal package
in cui il test era stato collocato.

La copertura della suite risultante (Prompt 2 e 3 combinati), misurata
con JaCoCo, risulta pari al 97.1\% statement e al 95.5\% branch — la
più alta tra le tre tecniche automatiche applicate a \texttt{Utf8}.

\paragraph{EvoSuite}

EvoSuite è stato eseguito su \texttt{Utf8} con la stessa
configurazione già descritta per \texttt{Resolver} (JDK 11, budget di
60 secondi), producendo 63 test. La \tabref{tab:evosuite-utf8}
riassume i risultati, e la \tabref{fig:jacoco-utf8-evosuite} ne mostra
il dettaglio per metodo.

La copertura, misurata isolatamente con JaCoCo, risulta pari al
95.2\% statement e al 95.5\% branch, sensibilmente più alta di quella
ottenuta su \texttt{Resolver} (88.5\%/84.2\%) con lo stesso budget di
ricerca. La differenza è coerente con quanto già osservato per le
altre tecniche automatiche: \texttt{Utf8}, essendo una classe più
piccola e con una logica meno ramificata rispetto a \texttt{Resolver},
risulta più facilmente esplorabile a parità di tempo di generazione.


\subsection{Integrazione tecnica di EvoSuite}
Integrare EvoSuite nel sistema di compilazione del progetto ha richiesto 
di superare diverse difficoltà tecniche. Di seguito descrivo in 
dettaglio ciascun problema incontrato e il metodo adottato per risolverlo.

\subsubsection{Compatibilità del bytecode nella fase di generazione}
La libreria interna a EvoSuite 1.2.0 che analizza il bytecode (chiamata ASM) 
non riesce a leggere correttamente le classi compilate con Java 21, 
che è il Java predefinito nel nostro ambiente di sviluppo. Inoltre, 
la libreria mostrava problemi con alcune dipendenze del progetto 
fornite come ``multi-release JAR'', ovvero archivi che contengono 
versioni differenti del codice per diverse edizioni di Java. 
Per ovviare a questo ostacolo, la fase di generazione dei test è stata modificata: 
il progetto viene compilato usando l'opzione di compatibilità per Java 11 
e per lanciare EvoSuite si adopera direttamente il programma binario di Java 11. 
Prima di procedere, inoltre, vengono eliminate dal percorso di ricerca 
delle classi (classpath) tutte le cartelle e i file denominati 
\texttt{META-INF/versions/*} presenti nei JAR multi-release che creavano conflitto.

\subsubsection{Gestione della dipendenza di esecuzione non disponibile}
Mentre per generare i test basta EvoSuite, per poterli effettivamente eseguire 
è indispensabile disporre della libreria \texttt{evosuite-standalone-runtime}. 
Purtroppo, la versione esatta di cui c'era bisogno non si trovava 
pubblicata sul repository centrale di Maven (Maven Central). 
Per risolvere questo problema senza compromettere il funzionamento automatico, 
è stato necessario inserire un passaggio specifico all'interno della pipeline di integrazione continua (CI). 
Questo comando scarica automaticamente il file JAR ufficiale direttamente 
dalla pagina delle versioni sul sito GitHub di EvoSuite e lo installa manualmente 
nel repository Maven locale della macchina di compilazione (\texttt{mvn install:install-file}) 
prima di iniziare la compilazione vera e propria. In questo modo, la build 
rimane perfettamente riproducibile anche partendo da un computer completamente pulito o da zero.

\subsubsection{Superamento delle restrizioni di sicurezza della macchina virtuale Java}
A partire dalla versione 18 di Java, l'utilizzo del comando interno 
\texttt{System.setSecurityManager(...)} — che il sistema di isolamento (sandbox) 
di EvoSuite sfrutta per funzionare in sicurezza — è stato bloccato per impostazione 
predefinita e richiede un permesso esplicito. 
Per aggirare questo blocco, si è reso necessario inserire l'apposito parametro 
di configurazione \texttt{-Djava.security.manager=allow} all'interno della riga di comando 
dei parametri (\texttt{argLine}) del plugin Surefire di Maven, che si occupa di eseguire i test. 
Visto che in Maven le configurazioni a valore singolo definite nel file \texttt{pom.xml} principale 
non si sommano ma vengono semplicemente sovrascritte da quelle del modulo figlio, 
questo flag di sicurezza è stato dichiarato direttamente ed esplicitamente 
nel file \texttt{pom.xml} del singolo modulo in cui i test vengono eseguiti.

\subsubsection{Verifica del bytecode durante l'esecuzione}
Quando i test generati da EvoSuite venivano eseguiti su Java 21, la macchina virtuale 
generava un errore specifico di controllo del codice noto come \texttt{VerifyError}, 
un problema che invece non compariva per niente se si utilizzava Java 11. 
Questo accade perché il codice modificato (chiamato \textit{instrumentato}) da EvoSuite 
non rispetta le regole di controllo e verifica molto più rigide introdotte nelle versioni di Java successive alla undicesima. 
La soluzione adottata è consistita nel configurare un elemento specifico chiamato \texttt{jdkToolchain} 
all'interno del plugin Surefire, imponendo al sistema di eseguire i test di EvoSuite 
sfruttando sempre e comunque Java 11, indipendentemente dalla versione di Java 
utilizzata per compilare il resto del software.

\subsubsection{Risoluzione dell'inquinamento di stato tra le classi di test}
L'introduzione dei test automatici di EvoSuite ha fatto emergere, in modo discontinuo 
e legato all'ordine con cui i test venivano eseguiti, un conflitto a livello del meccanismo \texttt{ServiceLoader} di Java. 
In pratica, il processo di modifica del codice di EvoSuite caricava alcuni moduli di configurazione SPI 
appoggiandosi a un gestore di classi differente rispetto a quello previsto dall'interfaccia originale. 
Questo finiva per corrompere i dati statici condivisi all'interno dell'applicazione 
(in particolare la classe \texttt{LogicalTypes}) per l'intera durata del processo di esecuzione di Surefire. 
Indagando a fondo, la causa è stata circoscritta ad alcuni file di risorse superflui 
situati nella cartella \texttt{META-INF/services}, rimasti nel progetto a causa di vecchi test nativi 
che erano già stati esclusi in precedenza (come descritto nella Sezione 2.3). 
La loro completa rimozione dalle risorse di test ha eliminato definitivamente la radice del problema.